\chapter{\IfLanguageName{dutch}{Stand van zaken}{State of the art}}%
\label{ch:stand-van-zaken}

% Tip: Begin elk hoofdstuk met een paragraaf inleiding die beschrijft hoe
% dit hoofdstuk past binnen het geheel van de bachelorproef. Geef in het
% bijzonder aan wat de link is met het vorige en volgende hoofdstuk.

% Pas na deze inleidende paragraaf komt de eerste sectiehoofding.
Bij mobiele applicaties voor buitendienstmedewerkers is het belangrijk dat de ingevoerde gegevens betrouwbaar zijn.
\textcite{KnowlesEtAl2014} onderzoeken dit bij medewerkers in het wegenonderhoud. \textcite{HardingEtAl2019} onderzoeken
hoe medewerkers in transportonderhoud gegevens gebruiken bij hun werk. Daarbij brengen ze de bestaande werkprocessen in kaart
via interviews en observaties. Deze onderzoeken gaan niet specifiek over HVAC-bedrijven, maar tonen wel waarom het belangrijk is om
eerst te begrijpen hoe medewerkers gegevens registreren en gebruiken. In deze bachelorproef wordt daarom eerst de huidige
werkwijze van het bedrijf onderzocht.

Om ook zonder netwerkverbinding gegevens te kunnen raadplegen, kan een mobiele applicatie gegevens lokaal opslaan.
De documentatie van Android beschrijft hoe een offline-first applicatie hiervoor een lokale gegevensbron combineert met
een gegevensbron op het netwerk. Wanneer gebruikers ook offline gegevens moeten kunnen aanpassen, kunnen wijzigingen
eerst lokaal worden bewaard en daarna in een wachtrij worden geplaatst. Zodra er opnieuw verbinding is, probeert de
applicatie de wijzigingen door te sturen~\autocite{Google2026OfflineFirst}. Ook \textcite{NingrumEtAl2026} bespreken
lokale opslag en uitgestelde synchronisatie als veelgebruikte onderdelen van applicaties met offlineondersteuning.
Bij offline gegevensregistratie is het belangrijk om wijzigingen blijvend op te slaan.
Alleen gegevens tijdelijk in het geheugen bewaren is niet
voldoende wanneer een applicatie wordt afgesloten. De documentatie van
Android beschrijft daarom lokale opslag op het apparaat en het bewaren
van werk dat later moet worden uitgevoerd~\autocite{Google2026OfflineFirst}.
Ook na verbindingsherstel kan een synchronisatiepoging mislukken.
De applicatie moet dan kunnen bepalen welke wijzigingen nog moeten
worden verstuurd. Dezelfde documentatie bespreekt hiervoor het opnieuw
proberen met steeds langere wachttijden tussen mislukte pogingen.
Dit voorkomt dat de applicatie onmiddellijk blijft proberen wanneer
de verbinding of server nog niet beschikbaar is~\autocite{Google2026OfflineFirst}.

Er bestaan verschillende manieren om voor synchronisatie in mobiele applicaties te zorgen.
\textcite{McCormickSchmidt2012} beschrijven onder meer wanneer gegevens worden gesynchroniseerd,
welke gegevens lokaal beschikbaar blijven en welke wijzigingen worden doorgestuurd. De geschikte
aanpak hangt vooral af van de netwerkverbinding en van hoe recent de gegevens moeten zijn.
Een wachtrij kan wijzigingen bewaren totdat ze kunnen worden verstuurd, maar lost niet automatisch conflicten op.
Een conflict kan bijvoorbeeld ontstaan wanneer twee medewerkers dezelfde gegevens aanpassen terwijl ze offline zijn.
Er zijn dan extra regels nodig om te bepalen welke wijzigingen worden behouden. In de documentatie van Android worden
onder andere het bijhouden van versies en last-write-wins besproken. Bij last-write-wins wordt één wijziging als de laatste
beschouwd en behouden, waardoor een andere wijziging kan worden overschreven~\autocite{Google2026OfflineFirst}.
Naast conflicten tussen gebruikers moet rekening worden gehouden met herhaalde verzending.
Een server kan een wijziging al hebben verwerkt terwijl het antwoord de mobiele applicatie niet bereikt.
Opnieuw versturen mag dan niet onbedoeld een tweede registratie veroorzaken. Een bewerking is idempotent wanneer meerdere
identieke verzoeken hetzelfde bedoelde effect hebben als één verzoek. Dit wordt beschreven in
de HTTP-specificatie~\autocite{FieldingEtAl2022}. Voor een materiaalregistratie is het bijvoorbeeld belangrijk om
onderscheid te maken tussen het opnieuw versturen van dezelfde registratie en het toevoegen van een nieuw verbruik.

Centrale conflictcontrole kan daarnaast voorkomen dat een wijziging op basis van verouderde gegevens een recentere wijziging overschrijft.
De HTTP-specificatie beschrijft hiervoor onder meer het gebruik van een ETag, een identificatie van een versie, samen met de
voorwaarde If-Match. De server controleert dan of de versie waarop het verzoek is gebaseerd nog overeenkomt met de huidige versie.
Als dat niet het geval is, kan het verzoek worden afgewezen. Deze controle detecteert een mogelijk conflict, maar bepaalt nog niet
hoe de verschillende wijzigingen inhoudelijk moeten worden afgehandeld~\autocite{FieldingEtAl2022}.

Een andere aanpak maakt gebruik van conflictvrije gerepliceerde datatypes, ook CRDT's genoemd. Daarmee kunnen verschillende
apparaten elk hun lokale kopie van gegevens aanpassen. Wanneer alle kopieën dezelfde wijzigingen hebben ontvangen, komen ze
onder de voorwaarden van het gekozen CRDT tot dezelfde toestand~\autocite{PreguicaEtAl2018}. Er bestaan verschillende CRDT's,
bijvoorbeeld voor tellers, verzamelingen en tekstgegevens. Elk type heeft regels die bepalen hoe wijzigingen worden samengevoegd.
De keuze van het datatype en de regels blijft dus belangrijk. Het automatisch samenvoegen van wijzigingen betekent niet dat
het resultaat ook altijd overeenkomt met wat de gebruikers bedoelden~\autocite{PreguicaEtAl2018}.

Bij CRDT's wordt onderscheid gemaakt tussen het uitwisselen van toestanden en het uitwisselen van bewerkingen. Bij de eerste aanpak
stuurt een apparaat zijn toestand door, waarna die met de lokale toestand van een ander apparaat wordt samengevoegd. Bij de tweede 
aanpak worden de uitgevoerde bewerkingen doorgestuurd en op de andere kopieën toegepast. Beide aanpakken stellen voorwaarden aan de
samenvoeging en de uitwisseling van gegevens~\autocite{PreguicaEtAl2018}.

Ook de betekenis van een bewerking speelt een rol. Het verhogen van een hoeveelheid is bijvoorbeeld een andere handeling dan het vervangen 
van die hoeveelheid door een nieuwe waarde. Een CRDT-teller kan verhogingen en verlagingen combineren, terwijl voor het gelijktijdig
vervangen van een waarde aanvullende keuzes nodig zijn~\autocite{PreguicaEtAl2018}.

\textcite{KleppmannEtAl2019} bespreken hoe lokale gegevensopslag en samenwerking tussen apparaten gecombineerd kunnen worden.
Ze noemen dit local-first software en onderzoeken hoe CRDT's hierbij kunnen helpen. Local-first gaat verder dan
alleen offline werken en houdt bijvoorbeeld ook rekening met de controle van gebruikers over hun gegevens.
Een wachtrij en CRDT's zijn daarbij geen volledig
aparte technieken. Ook een CRDT-gebaseerde applicatie moet wijzigingen bewaren en doorsturen. De vergelijking in dit onderzoek
gaat daarom over twee concrete implementaties die wijzigingen en conflicten op een andere manier
verwerken~\autocite{PreguicaEtAl2018}.

Overeenstemming tussen gegevenskopieën moet worden onderscheiden van het naleven van applicatieregels. \textcite{BailisEtAl2014} onderzoeken
onder welke voorwaarden bewerkingen zonder onderlinge afstemming kunnen worden uitgevoerd terwijl vastgelegde regels behouden blijven.
Hun analyse laat zien dat dit afhangt van de combinatie van de regels en de toegelaten bewerkingen. Sommige combinaties kunnen zonder
afstemming correct blijven, terwijl andere wel afstemming vereisen.

Toegepast op werkbonnen betekent dit dat automatisch samengevoegde gegevens afzonderlijk aan de bedrijfsvereisten moeten worden getoetst.
Bijvoorbeeld dat elke materiaalregistratie aan een bestaande werkbon gekoppeld blijft. Welke regels daadwerkelijk
gelden bij Goossens Warmte- en Koeltechniek, wordt tijdens de
bedrijfsanalyse vastgesteld.

De literatuur beschrijft verschillende mogelijkheden, maar geeft nog geen antwoord op welke aanpak het meest geschikt is
voor de werkbonnen van het onderzochte bedrijf. Dat apparaten na synchronisatie dezelfde gegevens bevatten, is bovendien niet
voldoende. Het resultaat moet ook voldoen aan de afspraken over werknotities en materiaalregistraties. Daarom worden de
bedrijfsvereisten gebruikt om te bepalen wat als een correct resultaat wordt beschouwd. Beide implementaties worden vervolgens
getest met dezelfde gegevens en situaties. Zo wordt onderzocht hoe ze verschillen in gegevenscorrectheid, synchronisatietijd
en implementatiecomplexiteit.

% \begin{figure}
%   \centering
%   \includegraphics[width=0.8\textwidth]{grail.jpg}
%   \caption[Voorbeeld figuur.]{\label{fig:grail}Voorbeeld van invoegen van een figuur. Zorg altijd voor een uitgebreid bijschrift dat de figuur volledig beschrijft zonder in de tekst te moeten gaan zoeken. Vergeet ook je bronvermelding niet!}
% \end{figure}

% \begin{listing}
%   \begin{minted}{python}
%     import pandas as pd
%     import seaborn as sns 

%     penguins = sns.load_dataset('penguins')
%     sns.relplot(data=penguins, x="flipper_length_mm", y="bill_length_mm", hue="species")
%   \end{minted}
%   \caption[Voorbeeld codefragment]{Voorbeeld van het invoegen van een codefragment.}
% \end{listing}

% \lipsum[7-20]

% \begin{table}
%   \centering
%   \begin{tabular}{lcr}
%     \toprule
%     \textbf{Kolom 1} & \textbf{Kolom 2} & \textbf{Kolom 3} \\
%     $\alpha$         & $\beta$          & $\gamma$         \\
%     \midrule
%     A                & 10.230           & a                \\
%     B                & 45.678           & b                \\
%     C                & 99.987           & c                \\
%     \bottomrule
%   \end{tabular}
%   \caption[Voorbeeld tabel]{\label{tab:example}Voorbeeld van een tabel.}
% \end{table}

